\subsection{Natural-image pixel-ridge landscape experiment}
\label{sec:app-vanhateren-landscape-methods}

This subsection describes the experiment underlying Fig.~\ref{fig:pixel} and
Figs.~\ref{fig:supp-pixel-de-mc}--\ref{fig:supp-pixel-selection}.  Its
Monte Carlo component uses actual natural-image patches, rather than Gaussian
designs having the same covariance.

\paragraph{Images, teacher, and population reference.}
Images, pools, covariance and disk teacher are those of
App.~\ref{sec:app-methods-images}: a 2,000-patch training pool, a disjoint
20,000-patch population pool that defines $\Sigma$ and is used to evaluate
every fitted model, and a disk teacher with
$S=\bstar^\top\Sigma\bstar=1736.34$.

\paragraph{Noise--regularization grid and deterministic equivalents.}
We evaluated 114 response-noise levels, comprising zero and a logarithmic grid
of $\sigma^2/S\in[10^{-8},10]$, and 478 ridge penalties.  We use the
scikit-learn convention
\begin{equation}
 \widehat\beta_{\alpha,\sigma}
 =X^\top(XX^\top+\alpha I_n)^{-1}\bm y_\sigma,
 \qquad \alpha=n\lambda,
\end{equation}
with $\alpha\in[10^{-5},10^7]$ ($\lambda\in[10^{-8},10^4]$).  The effective
penalty $\kappa$ shown in Fig.~\ref{fig:pixel} is the positive solution of
$n\lambda=\kappa\{n-\mathrm{df}_1(\kappa)\}$.  At each of the 54,492 grid
points we evaluated the deterministic equivalents using the eigenvalues of
the empirical population covariance and the disk teacher's coordinates in
that eigenbasis.  The prediction-optimal and control-optimal paths minimize,
respectively, the DE prediction error and the leading-DE accentuation error
over the finite penalty grid at each noise level.  All heatmaps and slices
are direct evaluations on this grid; no smoothing or interpolation was used.

\paragraph{Natural-image Monte Carlo.}
We ran 100 paired trials.  On trial $t$, using seed $20260918+t$, we sampled
$n=1,000$ patches without replacement from the 2,000-patch training pool,
centered every pixel using that sample, and drew a centered
$\epsilon\sim\mathcal N(0,I_n)$.  For every noise level we used the response
\begin{equation}
  \bm y_\sigma=X\bstar+\sigma\epsilon .
  \label{eq:app-vh-response}
\end{equation}
Thus the design and unit-noise direction were shared across the entire noise
axis within a trial.  This pairing reduces differences caused by resampling
and makes each trial a coherent path as $\sigma$ varies.

Because $n<d$, we diagonalized the $n\times n$ kernel matrix $XX^\top$ once
per trial.  For every penalty the fitted weight is exactly affine in the
noise scale,
\begin{equation}
 \widehat\beta_{\alpha,\sigma}
 =X^\top(XX^\top+\alpha I_n)^{-1}X\bstar
 +\sigma X^\top(XX^\top+\alpha I_n)^{-1}\epsilon
 =w_{0,\alpha}+\sigma w_{1,\alpha}.
 \label{eq:app-vh-affine-noise}
\end{equation}
We therefore cached the linear or quadratic coefficients in $\sigma$ of all
required quadratic forms and reconstructed the complete noise axis exactly.
This avoids refitting 54,492 models per trial and is not an interpolation.

Let $P\in\mathbb R^{20,000\times d}$ be the centered population-patch matrix,
$N=\bstar^\top\widehat\beta$, $D=\|\widehat\beta\|^2$,
$Q=\|P\widehat\beta\|^2/20{,}000$, and
$C=(P\bstar)^\top P\widehat\beta/20{,}000$.  We report
\begin{align}
 \frac{E_{gen}}S
 &=\frac{\|P(\widehat\beta-\bstar)\|^2}{20{,}000S},
 &R^2_{gen}&=1-\frac{E_{gen}}S,
 &\operatorname{slope}_{gen}&=\frac CQ,\\
 \frac{E_{acc}}S
 &=\left(1-\frac ND\right)^2,
 &R^2_{acc}&=1-\left(\frac DN-1\right)^2,
 &\operatorname{slope}_{acc}&=\frac ND.
 \label{eq:app-vh-metrics}
\end{align}
The fitted intercept is accounted for by sample centering and by the
cross-validation leverage below, but it is omitted from these population
weight-response metrics.

\paragraph{Empirical RidgeCV.}
For each trial and noise level we selected $\alpha$ by exact leave-one-out
cross-validation over the same 478-point grid.  If $XX^\top=U\operatorname{diag}(q_j)U^\top$,
the diagonal of the fitted-intercept ridge smoother is
\begin{equation}
 h_i(\alpha)=\sum_jU_{ij}^2\frac{q_j}{q_j+\alpha}+\frac1n,
\end{equation}
and the LOOCV loss is the mean of
$[(y_i-\widehat y_i)/(1-h_i)]^2$.  As in
Eq.~\ref{eq:app-vh-affine-noise}, this loss is quadratic in $\sigma$, so its
full landscape was reconstructed from three cached coefficients.  The green
path in Fig.~\ref{fig:pixel} is the median selected $\kappa$ across trials.
Metrics labeled empirical RidgeCV were instead evaluated at each trial's own
selected penalty before aggregation; they are not evaluations at the median
penalty.  We verified the selected penalties and minimum losses against
\texttt{sklearn.linear\_model.RidgeCV} at three noise levels.

\paragraph{Aggregation and checks.}
At every grid point we saved all 100 trial values and computed their
arithmetic mean, standard error, median, and 10th and 90th percentiles.  Bands
in Fig.~\ref{fig:pixel} show the 10--90\% trial range, not a confidence
interval; markers show the arithmetic mean.  We also checked the cached
noise-polynomial evaluation against an independent direct ridge solve for
all reported metrics.  Computation used float64 linear algebra on a GPU, and
plot-ready summaries were cached separately from the rendered figure.
